\bookchapter{Capítulo 69: Manual de consulta rápida e recursos}

\chapterintro{Este capítulo é um manual de consulta rápida que funciona por conta própria. Quando você não lembrar de uma sintaxe ou precisar tomar uma decisão rápida, venha direto para cá. Trate-o como um “cartão de bolso”: consulta de código frequente, perguntas frequentes e recursos de estudo, sem pular nenhuma página.}

\booksection{Cartão de consulta rápida de código frequente}

\begin{center}
\begin{tabularx}{\linewidth}{|X|X|}
\hline
\textbf{O que fazer} & \textbf{Código-chave} \tabularnewline
\hline
Estado privado da tela & @State private var x = valor inicial \tabularnewline
Lógica entre telas & \texttt{@Observable class VM \{ var x = 0 \}} \tabularnewline
Layout vertical/horizontal & \texttt{VStack \{ \}}/\texttt{HStack \{ \}} \tabularnewline
Espaço elástico & \texttt{Spacer()} \tabularnewline
Lista & List \{ ForEach(array) \{ ... \} \} \tabularnewline
Navegação entre telas & \texttt{NavigationStack + NavigationLink} \tabularnewline
Pop-up inferior (sheet) & .sheet(isPresented: \$show) \{ conteúdo \} \tabularnewline
Fechar a tela atual & \texttt{@Environment(\textbackslash{}.\allowbreak dismiss) var dismiss} \tabularnewline
Buscar dados pela rede & \texttt{try await URLSession.\allowbreak shared.\allowbreak data(from: url)} \tabularnewline
Interpretação de JSON & \texttt{struct X: Codable}+\texttt{JSONDecoder().\allowbreak decode} \tabularnewline
Salvar configurações pequenas & UserDefaults.standard.set(valor, forKey: "chave") \tabularnewline
Salvar dados estruturados & \texttt{@Model}+\texttt{modelContainer}+\texttt{@Query} \tabularnewline
Adaptação ao modo escuro & Use cores do sistema (.primary / Color(.systemGray6)), sem fixar valores \tabularnewline
Rodar testes & CmdU \tabularnewline
\hline
\end{tabularx}
\end{center}

\begin{center}\small\color{gray}Cartão de consulta rápida de desenvolvimento iOS (edição de agosto de 2026)\end{center}

\booksection{Perguntas frequentes (FAQ)}

P1: Não tenho Mac; consigo aprender desenvolvimento iOS?

Dá para aprender a linguagem, mas não para fazer o desenvolvimento completo. O Swift Playgrounds no iPad/iPhone cobre todo o conteúdo do Capítulo 61; porém o Xcode é exclusivo do macOS e, no fim, você vai precisar de um Mac. Com orçamento apertado, comece alugando um Mac na nuvem ou compre um Mac mini usado como transição, e decida o investimento depois de aprender a linguagem.

P2: Quanto custa aprender desenvolvimento iOS?

Fase de aprendizado: zero (Xcode gratuito + Apple ID gratuito para testar em aparelho real). Fase de publicação: conta de desenvolvedor a 99 dólares/ano (cerca de 700 yuans). Somando o hardware de um Mac, esse é o custo total. Diante do retorno — na App Store até uma ferramentinha tem usuários pagantes — a conta compensa.

P3: SwiftUI ou UIKit, qual aprender primeiro?

SwiftUI. Em 2026, novos projetos, novos tutoriais e novas vagas giram em torno do SwiftUI, e este livro o usa do começo ao fim. O UIKit é o framework da geração anterior e ainda vive em muitos projetos antigos — quando dominar o SwiftUI, complete com UIKit conforme a necessidade (há tutoriais oficiais de interoperabilidade); a curva será rápida. Não inverta a ordem.

P4: É fácil encontrar trabalho com desenvolvimento iOS?

O cenário de 2026: as vagas de início de carreira que exigem apenas “saber desenhar telas” são bem concorridas, mas desenvolvedores com “Swift + SwiftUI no domínio + um app publicado” continuam sendo cobiçados. O salário mediano das vagas iOS fica entre os primeiros das várias stacks há anos. Para quem está aprendendo, o melhor currículo é aquele App feito por você mesmo, na loja.

P5: Meu App foi rejeitado na revisão; e agora?

Fique calmo, é normal — a chance de rejeição na primeira submissão não é baixa. Os motivos da Apple costumam ser bem específicos (qual tela, qual diretriz); corrija item a item, responda com as explicações no App Store Connect e reenvie. Os três motivos de rejeição mais comuns: finalidade das permissões mal explicada, bugs na interface e metadados (capturas/descrição) que não correspondem ao produto. Checar esses três pontos com antecedência poupa uma rodada de espera.

\booksection{Atalhos de teclado mais usados do Xcode}

\begin{center}
\begin{tabularx}{\linewidth}{|X|X|}
\hline
\textbf{Ação} & \textbf{Atalho} \tabularnewline
\hline
Executar / testar & CmdR / CmdU \tabularnewline
Abrir o preview do canvas & OptionCmd \tabularnewline
Abrir arquivo rapidamente & CmdO \tabularnewline
Busca no projeto todo & CmdF \tabularnewline
Corrigir a indentação & CtrlI \tabularnewline
Comentar / descomentar & Cmd/ \tabularnewline
Limpar o cache de build & CmdK \tabularnewline
Ir para a definição & Cmd + clique \tabularnewline
Renomear (mudança de nome segura no projeto todo) & Clique direito Refactor Rename \tabularnewline
\hline
\end{tabularx}
\end{center}

\begin{center}\small\color{gray}Atalhos de teclado mais usados do Xcode\end{center}

\begin{tcbinfo}{Dica: erro "misterioso" de compilação? Limpe o cache primeiro}
Aquele fenômeno fantasma de “ontem compilava, hoje não compila mais” é, em seis ou sete casos de cada dez, cache de build corrompido. CmdK para limpar e compilar de novo resolve a maioria desses mistérios. Se não resolver, reinicie o Xcode (e o Mac junto) e, se persistir, “apague a pasta DerivedData” — a sequência final lendária passada de desenvolvedor iOS para desenvolvedor iOS.

\end{tcbinfo}

\booksection{Recursos de estudo recomendados}

\booksection{Recursos de estudo recomendados}

\begin{center}
\begin{tabularx}{\linewidth}{|X|X|X|}
\hline
\textbf{Tipo} & \textbf{Recurso} & \textbf{Por que recomendamos} \tabularnewline
\hline
Oficial & developer.apple.com & Documentação oficial da Apple; os tutoriais de SwiftUI/Swift estão todos lá, a maior autoridade \tabularnewline
Oficial & SwiftUI Tutorials da Apple & Curso oficial guiado passo a passo, nota máxima em facilidade para iniciantes \tabularnewline
Oficial & Site oficial do Swift (swift.org) & Documentação oficial da linguagem e playground online \tabularnewline
Comunidade & Hacking with Swift & Biblioteca de tutoriais gratuitos do Paul Hudson, qualidade reconhecidamente de primeira linha \tabularnewline
Comunidade & Stack Overflow (tag swiftui) & Primeiro porto de socorro para erros \tabularnewline
Vídeo & Canal oficial Apple Developer (vídeos da WWDC) & A explicação autoritativa das novidades de cada ano \tabularnewline
Design & Apple Human Interface Guidelines (HIG) & A bíblia oficial para criar interfaces com "jeitinho Apple" \tabularnewline
Complemento & Parte 7, "Prática de Programação com IA" (Capítulos 37-47) & O guia completo para desenvolver iOS com a ajuda da IA \tabularnewline
\hline
\end{tabularx}
\end{center}

\begin{center}\small\color{gray}Lista de recursos de estudo para desenvolvimento iOS (agosto de 2026)\end{center}

\booksection{Palavras finais: do primeiro App ao enésimo App}

Olhando para trás, você saiu de “não fazer ideia do que é SwiftUI” para instalar o Xcode, aprender Swift, desenhar telas com SwiftUI, conectar o App a um banco de dados e à rede, e ainda testar, empacotar e publicar — nesse caminho, você já venceu o trecho mais difícil da jornada do zero ao um.

O mundo da tecnologia muda rápido: frameworks evoluem, ferramentas se renovam e a IA assume cada vez mais o trabalho de “digitar”. Mas uma coisa não fica obsoleta: a sua capacidade de pegar uma ideia vaga, decompô-la em passos claros e realizá-la um a um. A ferramenta é o amplificador; o raciocínio é o alicerce. O amplificador pode ser trocado; o alicerce não pode balançar.

Você não é mais só "alguém que sabe usar celular". Você pode transformar as ideias da sua cabeça em apps de verdade rodando em centenas de milhões de iPhones. A partir de hoje, crie com o teclado — essa é a poesia de ser desenvolvedor iOS.

\booksection{Referências}

\begin{enumerate}
  \item Apple Inc. (2026). SwiftUI Framework Documentation. developer.apple.com.
  \item Apple Inc. (2026). Swift Programming Language Guide. swift.org.
  \item Apple Inc. (2026). SwiftData Framework Documentation. developer.apple.com.
  \item Apple Inc. (2026). App Store Review Guidelines. developer.apple.com/app-store/review.
  \item Apple Inc. (2026). Human Interface Guidelines (HIG). developer.apple.com/design.
  \item Apple Inc. (2026). Xcode 27 Release Notes and App Distribution Guide.
\end{enumerate}
